\section{Where the cost is}
\label{sec:problem}

The bound $2^{O(n)}\cdot M^{O(n^{2})}$ has two factors. We separate them by
counting, for a run, the number of distinct \emph{discrete shapes} --- nodes
taken up to their zone and their edge markers --- against the number of nodes.

\begin{table}[t]
\centering
\caption{Which factor is multiplying. The first row is a conjunction of $16$
independent bounded eventualities, the others nested $\AG(\EF\cdots)$ towers.}
\label{tab:regimes}
\begin{tabular}{@{}lrrr@{}}
\toprule
instance & shapes & nodes & zones per shape \\
\midrule
$16$ conjoined $\EF_{[a_j,b_j]}\,p_j$      & $262\,204$ & $263\,162$ & $\mathbf{1.00}$ \\
depth-$7$ $\AG(\EF_{[0,2]}\cdots)$ tower   & $5\,841$   & $10\,785$  & $1.85$ \\
$4$-level tower, varied windows            & $2\,635$   & $17\,370$  & $6.59$ \\
\bottomrule
\end{tabular}
\end{table}

Table~\ref{tab:regimes} is the paper's starting point. On the most expensive
instance the zone contributes nothing: one zone per discrete shape, and the
$263\,162$ nodes are $263\,162$ distinct discrete states. Abstraction,
widening and inclusion all act on the third column, and on the row that costs
the most the third column is $1.00$.

\paragraph{What the discrete states are.} Instrumenting the pursuit vector on a
depth-$6$ tower gives at most $12$ simultaneous pursuits, exactly two per
nesting level, with no modality duplicated, and a histogram of $|P|$ over the
nodes of $5, 22, 77, 172, 323, 545, 824, 952, 903, 803, 485, 167, 25$:
binomial, peaked in the middle. That is the signature of a search ranging over
the \emph{subsets} of a set of independent obligations. For
$\bigwedge_{j<k}\EF_{[a_j,b_j]}\,p_j$ the reading is immediate --- each
conjunct is outstanding or witnessed, the conjuncts do not interact, and the
node records the whole vector of $k$ bits. Measured node counts follow:
$1\,247$, $16\,863$, $66\,179$, $262\,975$, $1\,049\,619$ for
$k=8,12,14,16,18$.

\paragraph{There is nothing to learn from.} The natural import from SAT and SMT
is conflict-driven search. Counting how nodes end says it has no purchase.

\begin{table}[t]
\centering
\caption{How nodes end. Conflict-driven techniques need the middle columns.}
\label{tab:conflicts}
\begin{tabular}{@{}lrrr@{}}
\toprule
instance & accepted & rejected & closed \\
\midrule
$16$ conjoined eventualities (SAT)  & $262\,975$ & $\mathbf{0}$ & $\mathbf{0}$ \\
$4$-level tower (SAT)               & $12\,928$  & $1\,014$ & $790$ \\
depth-$6$ tower (SAT)               & $4\,281$   & $609$   & $385$ \\
$\AG(\EF_{[3,4]}p)\wedge\AG_{[4,5]}\neg p$ (UNSAT)
                                    & $161$      & $108$   & $53$ \\
\bottomrule
\end{tabular}
\end{table}

On the hardest instance every node is accepted and none is closed. The search
is not thrashing through contradictory subtrees that a learned clause would cut;
it is enumerating states that are all consistent. Asking a technique for
\emph{failing faster} to speed up a search that never fails is the wrong
request. What is needed is a technique for not enumerating a product --- and one
exists in the same literature, in model counting rather than in
solving~\cite{BacchusDalmaoPitassi2003,SangBeameKautz2005}.
